\documentclass[a4paper,11pt]{article}
\usepackage{exam-style-341}
\examinehead{341 农业综合知识三 · 网络技术分章练习}{第08章 应用层 参考答案与解析}
\begin{document}

\unitanswerscover{341 农业综合知识三 · 网络技术分章练习}{第08章 应用层 参考答案}{第08章练习的答案与解析；判断题均给出改正后的表述，综合题按邮件传送的实际环节给出完整过程。}

\answerpart{一、单项选择题（10 题，每题 3 分，共 30 分）}

\q A。\begin{analysis}DNS 默认使用 UDP 的 53 端口，以减少查询开销；当响应报文较长或进行区域传送时也可使用 TCP 的 53 端口（教材 6.1.1，印刷页 261—262）。B 25 是 SMTP；C 67 是 DHCP 服务器端口；D 80 是 HTTP。\end{analysis}

\q B。\begin{analysis}主机向本地域名服务器的查询一般采用递归查询；本地域名服务器向根域名服务器的查询通常采用迭代查询，即根服务器只告诉本地服务器下一步应找哪一台顶级域名服务器，由本地服务器继续查询（教材 6.1.3，印刷页 267）。A 递归查询用于主机与本地服务器之间；C 广播查询、D 反向查询不是本题所指的解析方式。\end{analysis}

\q B。\begin{analysis}FTP 服务器主进程打开熟知端口 21 等待客户连接，控制连接在整个会话期间保持打开（教材 6.2.2，印刷页 270）。A 20 是主动方式下服务器建立数据连接使用的端口；C 25 是 SMTP；D 80 是 HTTP。\end{analysis}

\q A。\begin{analysis}FTP 主动方式下，客户告诉服务器自己用于数据连接的端口，服务器用自己传送数据的熟知端口 20 与客户的端口建立数据连接（教材 6.2.2，印刷页 270）。B 21 是控制连接端口；C 22 是 SSH；D 23 是 TELNET。\end{analysis}

\q B。\begin{analysis}SMTP 用于发送（推送）邮件，默认端口 25（教材 6.5.2，印刷页 296）。A 21 是 FTP 控制连接；C 80 是 HTTP；D 110 是 POP3。\end{analysis}

\q C。\begin{analysis}POP3 用于从邮件服务器读取（拉取）邮件，默认端口 110（教材 6.5.4，印刷页 298）。A 25 是 SMTP；B 53 是 DNS；D 143 是 IMAP。\end{analysis}

\q A。\begin{analysis}MIME（通用互联网邮件扩充）规定了邮件正文的结构与编码方式，使电子邮件能够传送非 ASCII 文本、图像、音频、视频和附件（教材 6.5.6，印刷页 300）。B SMTP 只负责传送；C POP3 负责读取；D SNMP 是网络管理协议，与邮件格式无关。\end{analysis}

\q D。\begin{analysis}HTTP 默认端口为 80，HTTPS 默认端口为 443（教材 6.4.3，印刷页 276）。A 21 是 FTP 控制连接；B 25 是 SMTP；C 53 是 DNS。\end{analysis}

\q B。\begin{analysis}浏览器访问 Web 页面使用 HTTP（或加密的 HTTPS）协议，采用请求／响应方式获取网页、图片等资源（教材 6.4）。A FTP 用于文件传送；C SMTP 用于发送邮件；D DHCP 用于动态分配 IP 地址。\end{analysis}

\q C。\begin{analysis}DHCP 客户使用 UDP 端口 68，DHCP 服务器使用 UDP 端口 67（教材 6.6，印刷页 305）。A 53 是 DNS；B 67 是服务器端口；D 69 是 TFTP 的端口。\end{analysis}

\answerpart{二、多项选择题（3 题，每题 4 分，共 12 分）}

\q ABCD。\begin{analysis}域名服务器按层次分为根域名服务器、顶级域名服务器、权限域名服务器和本地域名服务器四类（教材 6.1.3，印刷页 265—266）。四项全对。\end{analysis}

\q ABC。\begin{analysis}SMTP 用于发送邮件，POP3 和 IMAP 用于读取邮件，三者都与电子邮件有关。D 错误：FTP 是文件传送协议，与电子邮件无关。\end{analysis}

\q ABC。\begin{analysis}HTTP 常用的请求方法有 GET、POST、HEAD、PUT、DELETE、OPTIONS 等（教材 6.4.3）。D 错误：SMTP 是邮件传送协议，不是 HTTP 请求方法。\end{analysis}

\answerpart{三、判断题（6 题，每题 2 分，共 12 分）}

\q ×。改正：DNS 默认使用 UDP 的 53 端口，但在响应报文较长或进行区域传送时也使用 TCP 的 53 端口，并非只能用 UDP。\begin{analysis}DNS 以 UDP 为主是为了减少开销，但 TCP 53 同样是 DNS 的标准端口（教材 6.1.1）。\end{analysis}

\q √。\begin{analysis}递归查询中，若本地域名服务器不知道被查询域名的 IP 地址，它以 DNS 客户的身份向其他服务器继续查询，最终把结果返回给主机，而不是让主机自己再查（教材 6.1.3，印刷页 267）。\end{analysis}

\q √。\begin{analysis}FTP 使用两条并行的 TCP 连接：控制连接传送命令与响应，数据连接传送文件内容，两者相互独立，控制信息带外传送（教材 6.2.2）。\end{analysis}

\q ×。改正：SMTP 是用于发送（推送）电子邮件的协议，接收邮件通常使用 POP3 或 IMAP。\begin{analysis}SMTP 负责把邮件从用户代理推送到邮件服务器、以及在邮件服务器之间传送；读取邮件由 POP3／IMAP 完成。\end{analysis}

\q √。\begin{analysis}MIME 在邮件首部增加内容类型、编码方式等字段，使邮件可以携带非 ASCII 文本、图像、音频和附件（教材 6.5.6）。\end{analysis}

\q √。\begin{analysis}HTTP 是无状态协议，服务器不保存请求之间的关联；Cookie 由服务器生成、浏览器保存并在后续请求中携带，从而识别用户、维持会话（教材 6.4.3，印刷页 282）。\end{analysis}

\answerpart{四、填空题（6 空，每空 1 分，共 6 分）}

\q （1）53。\begin{analysis}DNS 默认使用 53 端口（UDP，必要时 TCP）。\end{analysis}

\q （2）21。\begin{analysis}FTP 控制连接使用服务器端口 21。\end{analysis}

\q （3）25。\begin{analysis}SMTP 默认端口 25。\end{analysis}

\q （4）110。\begin{analysis}POP3 默认端口 110。\end{analysis}

\q （5）80。\begin{analysis}HTTP 默认端口 80（HTTPS 为 443）。\end{analysis}

\q （6）MIME。\begin{analysis}MIME 使电子邮件能够传送非 ASCII 数据与附件。\end{analysis}

\answerpart{五、简答题（2 题，每题 10 分，共 20 分）}

\q 参考答案要点：
（1）域名结构（3 分）：互联网采用层次树状结构的命名方法，域名由标号序列组成，标号之间用点分隔，从右到左级别降低。最右边是顶级域名（如 com、cn），其下是二级域名、三级域名……例如 www.hunau.edu.cn 中 cn 为顶级域名、edu 为二级域名、hunau 为三级域名。
（2）域名服务器分类（3 分）：根域名服务器（最高层次，知道所有顶级域名服务器的域名和 IP 地址）；顶级域名服务器（管理该顶级域名下的二级域名）；权限域名服务器（负责一个区）；本地域名服务器（主机首先查询的服务器，通常由 ISP 或单位设置）。
（3）解析过程（2 分）：主机向本地域名服务器查询；本地服务器若不能回答，则向根、顶级、权限域名服务器逐级查询，最后把结果返回主机，并在本地设置高速缓存以提高效率。
（4）递归查询与迭代查询（2 分）：递归查询由被查询的服务器代替发起者继续查询并返回最终结果（主机到本地服务器）；迭代查询由被查询的服务器返回“下一步应找哪台服务器”，由发起者自己继续查询（本地服务器到根、顶级、权限服务器）。

\q 参考答案要点：
（1）两条连接（4 分）：FTP 客户与服务器之间建立两条并行的 TCP 连接。控制连接使用服务器端口 21，在整个会话期间保持打开，用于传送 FTP 命令和响应，不传送文件内容；数据连接在每次传送文件时建立，用于传送文件内容，传送完毕后关闭。
（2）工作机制（3 分）：FTP 采用客户／服务器方式，服务器主进程在 21 端口监听，收到请求后启动从属进程处理；控制信息与数据分离（带外传送），使协议简单、便于在传输过程中进行控制。
（3）主动方式（2 分）：客户用 PORT 命令告诉服务器自己用于数据连接的端口，服务器用熟知端口 20 主动连接客户的该端口建立数据连接（教材 6.2.2 所述方式）。
（4）被动方式（1 分）：服务器打开一个随机端口并告知客户，由客户主动连接该端口建立数据连接；当客户位于防火墙之后时，被动方式更易穿透防火墙。（主动／被动为通行划分，教材重点讲述服务器端口 20 的主动方式。）

\answerpart{六、综合题（1 题，共 20 分）}

\q 参考答案与答题框架：

\textbf{【破题句】}该过程本质上是“用户代理—邮件服务器—邮件服务器—用户代理”的存储转发过程：发送用 SMTP（推送），接收用 POP3／IMAP（拉取），附件依赖 MIME 编码，各环节的服务器域名依赖 DNS 解析。

\textbf{（1）邮件传送环节、协议、方向与端口（8 分）}
\begin{analysis}
按四段给分，每段 2 分：
（1）A 的邮件客户端 $\rightarrow$ A 的发送方邮件服务器：使用 SMTP，推送方向，端口 25；
（2）A 的发送方邮件服务器 $\rightarrow$ B 的接收方邮件服务器：使用 SMTP，推送方向，端口 25（邮件服务器之间转发）；
（3）B 的接收方邮件服务器 $\rightarrow$ B 的手机邮件客户端：使用 POP3（端口 110）或 IMAP（端口 143），拉取方向；
（4）B 的手机邮件客户端通过上述连接收取邮件，邮件在服务器上以邮箱形式暂存（POP3 通常下载后删除服务器副本，IMAP 可在服务器上管理、多设备同步）。
\end{analysis}

\textbf{（2）MIME 与附件编码（4 分）}
\begin{analysis}
MIME 是通用互联网邮件扩充，它在邮件首部增加内容类型（Content-Type）、编码方式（Content-Transfer-Encoding）等字段，并规定邮件正文的分段结构（2 分）。传统 SMTP 只能传送 7 位 ASCII 文本，PDF 等二进制附件必须先用 Base64（或 quoted-printable）编码转换成 ASCII 字符，才能随邮件传送，接收方再解码还原（2 分）。
\end{analysis}

\textbf{（3）DNS 的作用（4 分）}
\begin{analysis}
邮件客户端和邮件服务器只知道对方邮件服务器的域名，发送前必须通过 DNS 把域名解析为 IP 地址（2 分）：先查询邮件服务器的 MX 记录确定负责收信的邮件服务器，再查询其 A 记录得到 IP 地址，然后才能建立 SMTP 连接；同理，B 的手机客户端也要通过 DNS 解析接收方邮件服务器的地址（2 分）。
\end{analysis}

\textbf{（4）Web 邮箱方式的区别（4 分）}
\begin{analysis}
使用浏览器登录 Web 邮箱时，用户代理是浏览器，用户与邮件服务器之间使用 HTTP／HTTPS（端口 80／443）提交和查看邮件，而不再使用 SMTP、POP3／IMAP（2 分）；SMTP 仍用于邮件服务器之间的转发，Web 服务器通过内部接口把邮件交给邮件服务器（2 分）。因此 Web 邮箱方式下“最后一公里”的协议由邮件专用协议变为 HTTP／HTTPS。
\end{analysis}

\textbf{【给分点】}（1）四段各 2 分；（2）MIME 作用 2 分、编码原因 2 分；（3）DNS 解析邮件服务器地址 2 分、MX／A 记录 2 分；（4）HTTP／HTTPS 替代专用协议 2 分、服务器之间仍用 SMTP 2 分。

\end{document}
